% Chapter 5 worked example. Included by ch05_resources.tex.
\section{Worked Example: Quality Signal After a Model and Workflow Update}
\label{sec:comp-ch05-worked-example}

\subsection{Purpose and Status}
This worked example demonstrates how the methods in Chapter~5 can be
translated into a reviewable institutional decision. All numerical inputs are
synthetic unless a source is identified explicitly. They are intended to show the
logic of the analysis, not to report empirical findings or establish a universal
threshold. Local use requires replacement of the assumptions with current,
versioned evidence and approval through the relevant clinical, managerial,
economic, legal, technical, and governance processes.

A deterioration-alert service has been used for eighteen months. Following an electronic health-record upgrade and a model update, the false-alert rate rises and nurses report more workarounds, while aggregate discrimination appears unchanged. 

\subsection{Decision Problem and Comparators}
\textbf{Decision problem.} Should the service continue unchanged, be restricted, undergo corrective action and revalidation, or be suspended pending investigation?

The comparator set is deliberately broader than the existing pathway. This prevents
a new technology from receiving a lower evidentiary threshold than staffing,
workflow, shared-service, conventional digital, or no-change alternatives.

\begin{longtable}{@{}p{0.08\textwidth}p{0.84\textwidth}@{}}
\caption{Comparators for the Chapter 5 worked example}
\label{tab:comp-ch05-comparators}\\
\toprule
\textbf{Option} & \textbf{Comparator or strategy} \\
\midrule
\endfirsthead
\toprule
\textbf{Option} & \textbf{Comparator or strategy} \\
\midrule
\endhead
\bottomrule
\endlastfoot
1 & Continue current operation \\
2 & Revert to the previous model and interface version \\
3 & Restrict use to selected units \\
4 & Suspend and return to the conventional escalation pathway \\
\end{longtable}

\subsection{Model and Decision Structure}
The analytical structure should follow the decision rather than the available
software. The team should map the causal pathway from intervention input to
information, human decision, action, outcome, resource consequence, and review.
The structure should also identify capacity constraints, uncertainty, subgroup
variation, implementation requirements, and events that would change the
recommendation.

A compact quantitative relationship used in this example is:
\begin{equation}
R_t=\frac{\text{events meeting the quality definition in period }t}{\text{eligible opportunities in period }t}
\label{eq:comp-ch05-key-relationship}
\end{equation}
The equation is an organizing aid rather than a substitute for disaggregated evidence,
qualitative findings, or mandatory safeguards. Its variables and interpretation must
be defined in the local protocol.

\subsection{Synthetic Parameters and Evidence Sources}
Table~\ref{tab:comp-ch05-parameters} provides the base-case inputs. A
production analysis should add parameter distributions, correlations, structural
alternatives, data dates, system versions, and evidence-confidence ratings.

\begin{longtable}{@{}p{0.32\textwidth}p{0.22\textwidth}p{0.38\textwidth}@{}}
\caption{Synthetic parameters for the Chapter 5 worked example}
\label{tab:comp-ch05-parameters}\\
\toprule
\textbf{Parameter} & \textbf{Base value} & \textbf{Source or interpretation} \\
\midrule
\endfirsthead
\toprule
\textbf{Parameter} & \textbf{Base value} & \textbf{Source or interpretation} \\
\midrule
\endhead
\bottomrule
\endlastfoot
Baseline false-alert rate & 14\% & Quality dashboard \\
Post-update false-alert rate & 23\% & Four-week monitoring window \\
Aggregate AUROC before/after & 0.82 / 0.81 & Monitoring report \\
Median response time before/after & 12 / 19 minutes & Workflow log \\
Units reporting workarounds & 5 of 8 & Safety walk-rounds \\
Serious incidents & 0 & Incident register \\
Near-miss reports & 7 & Patient-safety system \\
\end{longtable}

\subsection{Step-by-Step Analysis}
The sequence below is intended to make the analysis reproducible and to ensure
that evidence, economics, governance, and implementation develop together.

\begin{longtable}{@{}p{0.07\textwidth}p{0.55\textwidth}p{0.30\textwidth}@{}}
\caption{Analytical sequence}
\label{tab:comp-ch05-analysis-steps}\\
\toprule
\textbf{Step} & \textbf{Required work} & \textbf{Decision record} \\
\midrule
\endfirsthead
\toprule
\textbf{Step} & \textbf{Required work} & \textbf{Decision record} \\
\midrule
\endhead
\bottomrule
\endlastfoot
1 & Reconstruct the change timeline, data sources, model, software, interface, threshold, and workflow versions. & Evidence and responsible owner to be recorded \\
2 & Distinguish statistical stability from clinical and managerial acceptability. & Evidence and responsible owner to be recorded \\
3 & Investigate source-data, interface, threshold, user, staffing, and downstream-capacity explanations before selecting corrective action. & Evidence and responsible owner to be recorded \\
4 & Use the existing quality and patient-safety system to document containment, root causes, corrective and preventive action, and effectiveness checks. & Evidence and responsible owner to be recorded \\
5 & Assess whether the change is material and which validation evidence must be repeated. & Evidence and responsible owner to be recorded \\
6 & Communicate the decision, residual uncertainty, and conditions for return to unrestricted use. & Evidence and responsible owner to be recorded \\
\end{longtable}

\subsection{Illustrative Outputs}
The outputs in Table~\ref{tab:comp-ch05-outputs} show how technical,
clinical, operational, economic, equity, and governance findings should be kept
visible rather than compressed into a single favourable or unfavourable score.

\begin{longtable}{@{}p{0.30\textwidth}p{0.24\textwidth}p{0.38\textwidth}@{}}
\caption{Illustrative model and decision outputs}
\label{tab:comp-ch05-outputs}\\
\toprule
\textbf{Output} & \textbf{Result} & \textbf{Interpretation} \\
\midrule
\endfirsthead
\toprule
\textbf{Output} & \textbf{Result} & \textbf{Interpretation} \\
\midrule
\endhead
\bottomrule
\endlastfoot
Technical discrimination & Broadly stable & Does not rule out workflow or calibration failure \\
Operational performance & Deteriorated & Response delay and alert burden increased \\
Likely cause & Combined interface and threshold effects & Requires confirmation \\
Immediate action & Restrict and revert where feasible & Containment while investigation proceeds \\
Quality decision & CAPA and targeted revalidation & Not automatic retraining \\
\end{longtable}

\subsection{Sensitivity, Uncertainty, and Subgroups}
At minimum, the completed analysis should vary the parameters most likely to
change the decision, test at least one credible alternative model structure, and
report subgroup results separately where performance, access, response, burden,
or opportunity cost may differ. Deterministic analysis should identify thresholds;
probabilistic analysis should be used where credible parameter distributions and
correlations are available. Scenario analysis is preferable when structural or
future uncertainty cannot be represented defensibly by one probability model.

The record should distinguish uncertainty that can be reduced through evidence,
uncertainty that can be managed through safeguards, and uncertainty that remains
too consequential to accept. Missing evidence should not be represented as an
average or neutral value without justification.

\subsection{Interpretation and Recommendation}
Restrict use, restore the previous validated configuration where safe, and complete a formal incident and corrective-action process. Unrestricted use should resume only after the source of the signal is understood and the effectiveness of corrective action is verified.

The recommendation is conditional rather than permanent. It applies only to the
specified population, setting, intervention configuration, evidence cut-off, and
version. The following conditions should be incorporated into the approval,
contract, implementation, monitoring, or renewal record:

\begin{itemize}
 \item versioned reconstruction of data, software, model, interface, and workflow.
 \item unit-level and subgroup analysis.
 \item documented CAPA with named owners and completion dates.
 \item verification that response time and alert burden return within acceptable limits.
 \item executive authority to suspend or retire the intervention if performance remains unacceptable.
\end{itemize}

\subsection{Decision Statement}
\begin{quote}
For the specified decision problem and comparator set, the institution should
\emph{[proceed / proceed with conditions / redesign / pilot / restrict / defer /
replace / retire]}. The decision applies to \emph{[population, setting, version,
and evidence period]}, is owned by \emph{[accountable authority]}, and will be
reviewed on \emph{[date]} or earlier if \emph{[trigger events]} occur. The
evidence, assumptions, unresolved uncertainty, mandatory safeguards, monitoring
thresholds, and exit arrangements are recorded in the accompanying dossier.
\end{quote}
